\documentclass[11pt]{article}
\usepackage{amsmath,amsfonts}
\usepackage{graphicx}
%\usepackage{xcolor}
%\definecolor{gray}{rgb}{.75,.75,.75}
%\usepackage[landscape]{geometry}
%\usepackage{hyperref}
%\usepackage[bookmarks=true, pdfborder={0 0 .5}, urlbordercolor=gray]{hyperref}

\renewcommand{\thefootnote}{\fnsymbol{footnote}}


\addtolength{\topmargin}{-.9in}
\addtolength{\textheight}{1.6in}
\addtolength{\oddsidemargin}{-.4in}
\addtolength{\textwidth}{.8in}
\pagestyle{empty}

\newcommand{\ds}{\displaystyle}
\newcommand{\ts}{\textstyle}

\newcommand{\Z}{\mathbb{Z}}
\newcommand{\R}{\mathbb{R}}
\newcommand{\C}{\mathbb{C}}
\newcommand{\EE}{\mathcal{E}}
\newcommand{\tZ}{\tilde{\Z}}

\newcommand{\sech}{\mathrm{sech}}
\newcommand{\csch}{\mathrm{csch}}


\newcommand{\Sym}{\mathrm{Sym}}

\renewcommand{\circle}[1]{{\bigcirc\hspace{-.75em}#1}}

\begin{document}

\footnote[0]{Notes by Peter Magyar \ {\tt magyar@math.msu.edu}}
\vspace{-1em}

\centerline{\bf Math 133\hspace{1.5em} \hfill{\Large\bf
Ratio Test}\hfill Stewart \S11.6/I}
\vspace{1em}

\noindent 
We have one more important test for convergence of 
an infinite series $\sum_{n=1}^\infty a_n$. 
This test does not require us to 
choose a comparison series: instead, we test the ratio of 
each term $a_n$ compared to the next term $a_{n+1}$.
\\[-1em]
\begin{quote}
{\it Ratio Test:} Suppose $\ds\lim\limits_{n\to\infty} \left|\frac{a_{n+1}}{a_n}\right|=L$\,.
\\[-1em]
\begin{itemize}
\item If $L<1$, then $\sum_{n=1}^\infty a_n$ converges.
\item If $L>1$, then $\sum_{n=1}^\infty a_n$ diverges.
\item If $L=1$, then this test fails to determine convergence.
\end{itemize}

\end{quote}

\noindent
{\it Proof:} Assuming $a_n>0$, the limit $\lim\limits_{n\to\infty} \left|\frac{a_{n+1}}{a_n}\right|=L$ means that, for any small number $\epsilon>0$,
we can take a starting point $N$ so that  for all $n\geq N$, we have:
$$\begin{array}{ccccc}
L{-}\epsilon & \leq\ & \dfrac{a_{n+1}}{a_n}& \leq\ &  L{+}\epsilon 
\\[1em]
a_n(L{-}\epsilon) & \leq\ &  a_{n+1}& \leq\ &  a_n(L{+}\epsilon)\,.
\end{array}$$
Iterating this inequality gives: 
$
c_1(L{-}\epsilon)^n  \leq\  a_{n} \leq\   c_2(L{+}\epsilon)^n
$
for some constants $c_1,c_2$.\footnote{
Specifically:
$
a_{n}\leq a_{n-1}(L{+}\epsilon)\leq a_{n-2}(L{+}\epsilon)^2
\leq\cdots\leq a_N(L{+}\epsilon)^{n-N} = \frac{a_N}{(L{+}\epsilon)^{N}} (L{+}\epsilon)^n.
$
}

If $L<1$, we take $\epsilon$ small enough
that $L{+}\epsilon<1$, and we compare $\sum a_n$ to the convergent ceiling
series $\sum c_2(L{+}\epsilon)^n$. 
If $L>1$, we take $\epsilon$ small enough
that $L{-}\epsilon>1$, and we compare $\sum a_n$ to the divergent floor
series $\sum c_2(L{-}\epsilon)^n$. 
If $L=1$, adding any $\epsilon$ produces a divergent ceiling, and subtracting any $\epsilon$ produces a convergent floor, neither of which would constrain the original  series.
Finally, for the general case where the $a_n$'s may be positive or negative, 
the above argument shows $\sum |a_n|$ converges, which implies $\sum a_n$ converges by \S11.6 Part II. Q.E.D.
 \\[.9em]
 The Ratio Test is most useful when $a_n$ is a product of a growing number of factors, which will mostly cancel out in $\frac{a_{n+1}}{a_n}$.
 \\[.2em]
 {\sc example:} Determine the convergence of $\ds\sum_{n=1}^\infty \frac{n^2}{2^n}$\,.
 
We did this one in \S11.4 by finding a tricky comparison series. The Ratio Test naturally applies here,
 because $a_n= \frac{n^2}{2^n}=(n)(n)(\frac12)\cdots(\frac12)$ has more and more factors as $n$ gets larger.
We have:
$$
L\ =\
\lim_{n\to\infty} \left|\frac{a_{n+1}}{a_n}\right|
\ =\
\lim\limits_{n\to\infty} 
\frac{(n{+}1)^2}{2^{n+1}}\left/
\frac{n^2}{2^n}
\right.
\ =\
\lim_{n\to\infty}
\frac{(n{+}1)^2}{n^2}\left/
\frac{2^{n+1}}{2^n}
\right.
\ =\
\frac12\,.
$$
Since $L=\frac12<1$, the Test shows $\sum a_n$ converges.
\\[.5em]
 {\sc example:} Determine the convergence of $\ds\sum_{n=1}^\infty (-1)^n\frac{x^{2n}}{n!}$, where $x$ is a given number and 
 we use the factorial notation $n! = (n)(n{-}1)(n{-}2)\cdots (2)(1)$. 
Again, the terms have a large number of factors, so we use the Ratio Test:
$$
L\ =\
\lim_{n\to\infty} \left|\frac{a_{n+1}}{a_n}\right|
\ =\
\lim\limits_{n\to\infty} 
\frac{x^{2(n+1)}}{(n{+}1)!}\left/
\frac{x^{2n}}{n!}
\right.
\ =\
\lim\limits_{n\to\infty} \frac{x^2}{n{+}1}
\ =\ 0\,.
$$
Since $L = 0<1$, the Test shows $\sum a_n$ converges.

\end{document}

%%%%%%%%%%%  Picture  %%%%%%%%%%
\\[0em]
\centerline{
%\includegraphics[width=2in]
{1-8-Discont-iv.png}
\qquad \qquad 
%\includegraphics[width=2in]
{1-8-Discont-v.png} 
}
\vspace{0em}

\noindent
%%%%%%%%%%%%%%%%%%%%%%%%%

%%%%%%%%%%%  Table  %%%%%%%%%%
In fact, we get the following derivatives:
$$\begin{array}{|c||c|c|c|c|c|c|}
\hline &&&&&& \\[-.9em] 
f(x) & \sin(x) & \cos(x) & \tan(x)  & \sec(x) & \csc(x) & \cot(x)
\\[.2em] \hline &&&&&& \\[-.9em]
f'(x) &\cos(x) & -\sin(x) & \sec^2(x) & \tan(x)\sec(x) & -\cot(x)\csc(x) & -\csc^2(x)
\\[.1em] \hline
\end{array}$$
\\[1em]
%%%%%%%%%%%%%%%%%%%%%%%%%%



















